\section{Conclusion}
\label{sec:conclusion}

\subsection{The Argument, Restated}

The gerrymandering problem is not a problem of bad actors.
It is a problem of a system that gives actors with bad incentives---legislators
who draw their own districts---the authority to make consequential choices.
The solution is not to hope for better actors.
It is to remove the choices.

The \dia{} removes the choices by mandating an algorithm whose identity is
determined by first principles that both parties have already endorsed.

The four steps are simple.
Congressional districts must have equal population---\wesberry{} said so in 1964.
The correct method for apportioning equal-population units to geographic
sub-units in proportion to population is Huntington-Hill---Congress said so
in 1941 and the Supreme Court confirmed it in 1992.
Both parties agree that politicians should not draw their own districts---
they have said so, on the record, in courts across six decades.
The interstate and intrastate apportionment problems are structurally
isomorphic---logic says so, and the formal proof is in
Section~\ref{sec:why-apportionregions}.

Therefore: apply Huntington-Hill within states.
The algorithm that implements this is \ar{}.
The law that mandates it is the \dia{}.

\subsection{Why This Moment}

Three events in the past seven years have created the conditions for the
\dia{} to succeed.

\paragraph{Rucho, 2019.}
\rucho{} v.\ \textit{Common Cause}, 588 U.S. 684 (2019), closed the federal
courts to partisan gerrymandering claims.
This was not the end of redistricting litigation---it redirected plaintiffs
to state courts, which now manage fifty parallel redistricting jurisprudences.
But it sent a signal: the Supreme Court will not resolve the gerrymandering
problem through constitutional adjudication.
Congress must act.
\rucho{} created the legislative opening that the \dia{} fills.

\paragraph{Callais, 2026.}
\callais{}, 608 U.S. \_\_\_ (2026), established the disentanglement
requirement at page 36.
This requirement might have seemed to complicate algorithmic redistricting---
how can an algorithm satisfy both racial and partisan neutrality requirements
simultaneously?
Instead, it simplifies the problem by pointing to exactly the architecture
that \ar{} implements: a default run with zero signal of either type, and
a separate VRA overlay for states where the Gingles conditions are met.
\callais{} did not foreclose algorithmic redistricting.
It prescribed it.

\paragraph{The empirical record, 2026.}
The B-series research program \citep{b0paper,b7paper,b11paper,b13paper,b15paper}
has now produced the empirical foundation that prior reform proposals lacked.
\ar{} has zero seed variance.
Its solution space is certified as narrow.
Two-thirds of states produce census-stable maps.
The Wisconsin bakeoff proves that algorithm selection is a political choice,
and that \ar{}---the constitutionally derived algorithm---produces a more
proportional outcome in the most-litigated state in the country.
The empirical record is complete.

\subsection{The Objection That Will Be Raised}

The objection that will be raised is this: Huntington-Hill is not in the
Constitution.
The Constitution says ``as nearly as is practicable'' (\wesberry{}) and
``according to their respective Numbers'' (Article~I, \S2).
It does not say ``Huntington-Hill.''
So why is \ar{}, the Huntington-Hill implementation, constitutionally derived,
while GeoSection, the compactness-optimal implementation, is merely a
practitioner's choice?

The answer is that constitutional derivation does not require explicit
constitutional text for every implementation detail.
\hh{} is constitutionally derived because:
\begin{enumerate}
  \item Congress mandated it for the interstate problem in 2~U.S.C.~\usca{}~2a(a).
  \item The Supreme Court upheld it in \textit{Department of Commerce
    v.\ Montana}.
  \item No alternative has been enacted or constitutionally endorsed for
    the same problem.
\end{enumerate}
GeoSection is a practitioner's choice because none of these three predicates
is true of it: Congress has not mandated it, the Supreme Court has not
endorsed it, and there are several alternatives (AreaSection, unweighted
bisection, etc.) that are equally defensible on other grounds.

The constitutional derivation of \ar{} is not textual.
It is doctrinal: \hh{} is the method that the political branches have already
selected, through the ordinary democratic process, for the structurally identical
problem.
Applying it to the intrastate problem requires no new democratic selection.
It requires only recognizing that the same problem calls for the same solution.

\subsection{The Promise of the One-Sentence Law}

The history of redistricting reform is a history of complex proposals.
Independent commissions: How are they appointed? Who appoints them?
What criteria do they apply? What happens when commissioners disagree?
Mathematical optimization: Which objective function? Which constraints?
What weights? Which solver?

Every question in this history is an opportunity for political capture.
Every answer is a political choice.

The \dia{}'s one-sentence operative provision eliminates the questions.
There is no commission to appoint.
There is no objective function to choose.
There is no weight to set (beyond the county-preservation default, which is
itself a function of an existing administrative standard).
There is no solver to select: METIS is the industry standard, open-source,
and its behavior is fully documented.

The promise of the one-sentence law is not that it will produce perfect maps.
No algorithm produces perfect maps.
The promise is that it will produce maps that no one can blame on anyone.
The algorithm is what the census says it should be.
The boundaries are where the population lives.
The districts are as compact as geography permits.
The result is as proportional as the single-member-district system allows,
given the state's actual population distribution.

That is not everything.
But it is enough to stop the cycle of partisan map-drawing, partisan
litigation, and partisan map re-drawing that has consumed the American
redistricting process for two centuries.

\subsection{The Call to Action}

This Article calls on Congress to enact the \dia{} in the form proposed
in Section~\ref{sec:implementation}.
The operative text is:

\begin{statute}
\noindent\textbf{Proposed 2 U.S.C.\ \usca{}\ 2a(c) [new].}
Each State shall apportion its congressional districts by applying the
Huntington-Hill priority rule to its census tracts, using as input only
the prime factorization of the State's apportioned seat count and the
federal TIGER/Line adjacency graph for the census year in which the
apportionment was performed, with a seed determined by the formula
\texttt{SHA-256(census\_release\_id || `DIA\_SEED\_V1')}.
\end{statute}

This is not a partisan proposal.
It is the logical completion of the apportionment framework that Congress
began in 1941.
It uses the same method.
It solves the same problem.
It removes the same source of distortion.
It does so at the next level of the hierarchy.

The gerrymandering problem has a solution.
It is sixty-seven words long.
The time to enact it is now.
